\documentclass[12pt,a4paper]{article}
\usepackage[utf8]{inputenc}
\usepackage[T2A]{fontenc}
\usepackage[russian]{babel}
\usepackage{geometry}
\usepackage{graphicx}
\usepackage{amsmath}
\usepackage{amsfonts}
\DeclareMathOperator{\rank}{rank}
\usepackage{booktabs}
\usepackage{hyperref}
\hypersetup{
    pdftitle={PLA Lab 6 Control Questions Answers},
    pdfauthor={Anufriev Andrey Sergeevich},
    pdfsubject={Singular Value Decomposition, control questions},
    linkcolor=blue,
    filecolor=magenta,
    urlcolor=cyan
}
\usepackage{paracol}
\usepackage{src/title}

\geometry{top=2cm, left=2cm, right=2cm, bottom=2cm}

\title{Ответы на контрольные вопросы по лабораторной работе №6 (Сингулярное разложение матриц)}

\author{%
    Аннуфриев Андрей Сергеевич (465029)\\
    Бахуров Николай Алексеевич (504862)\\
    Раджапов Фёдор Сергеевич (502544)%
}
\supervisor{Зарина Ольга Андреевна}
\date{01.11.2026}

\abstract{Развёрнутые ответы на восемь контрольных вопросов лабораторной
работы №6 по практической линейной алгебре, посвящённой сингулярному
разложению матриц и его применениям (сжатие изображений, LSA).}

\keywords{сингулярное разложение; ортогональные матрицы; симметричные
матрицы; псевдообращение; LSA; образ и ядро}

\begin{document}
\sloppy
\maketitle

\section{Вопрос 1. Свойства ортогональных матриц}
Матрица $Q$ называется ортогональной, если $Q^\top Q = I$ (эквивалентно
$QQ^\top = I$ и $Q^{-1} = Q^\top$). Основные свойства:
\begin{itemize}
    \item столбцы (и строки) ортогональной матрицы образуют ортонормированный
    базис: попеременно ортогональны и имеют единичную длину;
    \item обратная матрица равна транспонированной: $Q^{-1}=Q^\top$;
    \item определитель равен $\pm 1$ ($\det Q = 1$ для поворотов, $-1$ для
    отражений);
    \item ортогональные преобразования сохраняют скалярное произведение,
    длины векторов и углы между ними (изометрии с неподвижным началом);
    \item собственные числа ортогональной матрицы по модулю равны 1;
    \item произведение ортогональных матриц ортогонально;
    \item матрицы группы $O(2)$ — это повороты ($SO(2)$, $\det=1$) и
    отражения ($\det=-1$).
\end{itemize}

\section{Вопрос 2. Свойства симметричных матриц; связь собственных чисел $AA^\top$ и $A^\top A$}
Симметричная матрица ($M=M^\top$) обладает следующими свойствами:
\begin{itemize}
    \item все собственные числа вещественны;
    \item собственные векторы, соответствующие различным собственным числам,
    ортогональны;
    \item матрица диагонализируется ортогонально (спектральная теорема):
    $M = Q\Lambda Q^\top$, где $Q$ ортогональна, $\Lambda$ диагональна;
    \item при существовании полного набора собственных векторов они образуют
    ортонормированный базис.
\end{itemize}
Связь: ненулевые собственные числа матриц $AA^\top$ и $A^\top A$
совпадают. Действительно, если $A^\top A v = \lambda v$ при $\lambda>0$,
то $AA^\top (Av) = \lambda(Av)$, и $Av\neq0$ — собственный вектор
$AA^\top$ с тем же $\lambda$. Аналогично в обратную сторону. нулевые
собственные числа могут отличаться по кратности (размерности ядер), если
$m\neq n$. Поскольку сингулярные числа $\sigma_i$ определяются как
$\sigma_i = \sqrt{\lambda_i(A^\top A)} = \sqrt{\lambda_i(AA^\top)}$, они
однообразны для обеих матриц.

\section{Вопрос 3. Геометрический смысл неквадратных матриц; ранг произведения}
Неквадратная матрица $A$ размера $m\times n$ задаёт линейное отображение
из $\mathbb R^n$ в $\mathbb R^m$ (из пространства с одной размерностью
в пространство с другой). Оно может:
\begin{itemize}
    \item вложить пространство меньшей размерности в большее ($m>n$),
    не используя всё целевое пространство (образ — подпространство
    размерности $\rank A \le n$);
    \item проецировать большее пространство в меньшее ($m<n$),
    сжимая некоторые направления в ноль (ядро нетривиально, если $n>m$);
    \item не быть обратимым ни слева, ни справа, если не квадратна.
\end{itemize}
Операции с образом/ядром: $\operatorname{Im}(A)\subseteq\mathbb R^m$ —
натянуто столбцами $A$; $\dim\operatorname{Im}(A)=\rank A$;
$\operatorname{Ker}(A)\subseteq\mathbb R^n$ — решения $Ax=0$;
$\dim\operatorname{Ker}(A)=n-\rank A$.

Ранг произведения: для матриц $A$ ($m\times p$) и $B$ ($p\times n$)
\[
    \rank(AB)\le\min\{\rank A,\rank B\},
\]
а более точно (неравенство Сильвестра):
\[
    \rank A + \rank B - p \le \rank(AB) \le \min\{\rank A,\rank B\}.
\]
Ранг произведения не может превышать ранга любого сомножителя; может быть
строго меньше (например, проекция на ядро первого отображения).

\section{Вопрос 4. Сингулярное разложение; геометрический смысл; смысл сингулярных чисел}
Сингулярное разложение матрицы $A$ ($m\times n$):
\[
    A = U\Sigma V^\top,
\]
где $U\in O(m)$, $V\in O(n)$, $\Sigma$ — $m\times n$ диагональная с
$\sigma_1\ge\dots\ge\sigma_r>0$, $r=\rank A$.

Геометрический смысл: отображение $x\mapsto Ax$ разлагается на три шага:
\begin{enumerate}
    \item $V^\top$ — поворот (или отражение) в исходном пространстве,
    переходящий к базовой системе правых сингулярных векторов;
    \item $\Sigma$ — независимое масштабирование по осям на коэффициенты
    $\sigma_i$ (растяжение/сжатие; $\sigma_i=0$ сжимает направление в ноль);
    \item $U$ — поворот (или отражение) в целевом пространстве, переводящий
    оси в направления левых сингулярных векторов.
\end{enumerate}
Сингулярные числа $\sigma_i$ показывают, во сколько раз отображение
растягивает единичный отрезок вдоль $i$-го собственного направления
(правого сингулярного вектора); величина $\sigma_1 = \|A\|_2$ — нормы
оператора.

\section{Вопрос 5. Связь SVD с Im (Range) и Ker (Nullspace)}
Пусть $A=U\Sigma V^\top$, $r=\rank A$, $\sigma_1,\dots,\sigma_r>0$.
\begin{itemize}
    \item $\operatorname{Im}(A) = \langle u_1,\dots,u_r\rangle$ — натянуто
    на первые $r$ левых сингулярных векторов;
    \item $\operatorname{Ker}(A) = \langle v_{r+1},\dots,v_n\rangle$ — натянуто
    на правые сингулярные векторы с нулевыми сингулярными числами;
    \item $\dim\operatorname{Im}(A)=r$, $\dim\operatorname{Ker}(A)=n-r$;
    \item $\operatorname{Im}(A^\top)=\langle v_1,\dots,v_r\rangle$,
    $\operatorname{Ker}(A^\top)=\langle u_{r+1},\dots,u_m\rangle$.
\end{itemize}
Таким образом, правые сингулярные векторы с $\sigma_i=0$ порождают ядро,
а левые — ортогональное дополнение к образу в целевом пространстве.

\section{Вопрос 6. Усечённое SVD; равенство с исходной матрицей; практическое применение}
Усечённое SVD ранга $k\le r$:
\[
    A_k = U_k\Sigma_k V_k^\top = \sum_{i=1}^{k}\sigma_i u_i v_i^\top.
\]
Результат совпадает с исходной матрицей тогда и только тогда, когда $k=r$
(количество ненулевых сингулярных чисел), то есть все ненулевые компоненты
учтены; при $k<r$ это наилучшее приближение ранга $k$ (теорема Эккарта —
Янгу) в нормах Фробениуса и спектральной.

Практическое применение:
\begin{itemize}
    \item сжатие данных и изображений: хранят $U_k,\Sigma_k,V_k^\top$
    ($k(m+n+1)$ чисел вместо $mn$);
    \item шумоподавление: отбражение мелких $\sigma_i$, соответствующих шуму;
    \item псевдообращение и решение неполноранговых систем;
    \item латентный семантический анализ, анализ главных компонент;
    \item вычисление ранга и числа обусловленности.
\end{itemize}

\section{Вопрос 7. Псевдообращение через SVD; свойства псевдообратных матриц}
Псевдообратная матрица Мура — Пенроуза $A^+$ размера $n\times m$ через
SVD $A=U\Sigma V^\top$:
\[
    A^+ = V\Sigma^+U^\top,
\]
где $(\Sigma^+)_{ii}=1/\sigma_i$ при $\sigma_i>0$, и $0$ при $\sigma_i=0$;
все остальные элементы $\Sigma^+$ — нули.

Свойства псевдообратных матриц (четыре условия Мура — Пенроуза):
\begin{itemize}
    \item $AA^+A = A$;
    \item $A^+AA^+ = A^+$;
    \item $(AA^+)^\top = AA^+$;
    \item $(A^+A)^\top = A^+A$.
\end{itemize}
Дополнительно: $AA^+$ — ортогональный проектор на $\operatorname{Im}(A)$;
$A^+A$ — ортогональный проектор на $\operatorname{Im}(A^\top)$
(т. е. на ортогональное дополнение к $\operatorname{Ker}(A)$);
$A^+$ даёт решение наименьшей нормы неполноранговой системы и решение
наименьших квадратов переопределённой системы.

\section{Вопрос 8. Смысл левых и правых сингулярных векторов в LSA; смысл сингулярных чисел}
В LSA терм-документная матрица $A$ (термины $\times$ документы)
разлагается как $A=U\Sigma V^\top$:
\begin{itemize}
    \item \textbf{Левые сингулярные векторы} (столбцы $U$) — ортонормированные
    направления в \emph{термин-пространстве}; каждый столбец $u_i$ задает
    латентную тему: компоненты $u_i$ показывают, какие слова характеризуют
    эту тему и с каким знаком/весом.
    \item \textbf{Правые сингулярные векторы} (столбцы $V$, строки $V^\top$)
    — координаты документов в \emph{документ-пространстве}; компоненты $v_i$
    показывают, насколько каждый документ выражает тему $i$.
    \item \textbf{Сингулярные числа} $\sigma_i$ — вес (значимость)
    соответствующей латентной темы в коллекции: чем больше $\sigma_i$, тем
    сильнее тема $i$ выражена в данных.
\end{itemize}
Поэтому первые два сингулярных числа и пары векторов $(u_1,v_1)$,
$(u_2,v_2)$ задают две доминирующие темы коллекции: топ-5 слов по $|u_i|$
определяют содержание темы, а топ-документов по $|v_i|$ — документы, наиболее
полностью соответствующие теме.

\end{document}
